\[ %%%%%%%%%%%%%%%%%%%%%%%%%%%%%%%%%%%%%%%%%%%%%%%%%%%%%%%%%%%%%%%%%%%%%%%%%%%%%%%% \subsection{Physical Consistency Evaluation} \label{sec:physical_consistency} %%%%%%%%%%%%%%%%%%%%%%%%%%%%%%%%%%%%%%%%%%%%%%%%%%%%%%%%%%%%%%%%%%%%%%%%%%%%%%%% Prediction accuracy alone is insufficient for evaluating scientific machine learning models. A physically plausible model must not only approximate observed trajectories, but also preserve the underlying constraints imposed by the governing equations. Therefore, this experiment evaluates whether PHYSGEN maintains physical consistency during long-horizon autonomous prediction. The evaluation criterion is: \begin{equation} \hat{X}_t \in \mathcal{M}_{phys}, \end{equation} where $\mathcal{M}_{phys}$ represents the admissible physical state manifold. %%%%%%%%%%%%%%%%%%%%%%%%%%%%%%%%%%%%%%%%%%%%%%%%%%%%%%%%%%%%%%%%%%%%%%%%%%%%%%%% \subsubsection{Physical Constraint Evaluation Framework} %%%%%%%%%%%%%%%%%%%%%%%%%%%%%%%%%%%%%%%%%%%%%%%%%%%%%%%%%%%%%%%%%%%%%%%%%%%%%%%% For each dynamical system, physical consistency is evaluated using the corresponding governing constraints. The general physical violation metric is defined as: \begin{equation} E_{phys} = \frac{1}{T} \sum_{t=1}^{T} \| \mathcal{F}(\hat{X}_t) \|^2, \end{equation} where: \begin{itemize} \item $\mathcal{F}$ represents the physical constraint operator; \item $\hat{X}_t$ is the predicted state. \end{itemize} Lower values indicate better physical consistency. %%%%%%%%%%%%%%%%%%%%%%%%%%%%%%%%%%%%%%%%%%%%%%%%%%%%%%%%%%%%%%%%%%%%%%%%%%%%%%%% \subsubsection{Conservation Law Preservation} %%%%%%%%%%%%%%%%%%%%%%%%%%%%%%%%%%%%%%%%%%%%%%%%%%%%%%%%%%%%%%%%%%%%%%%%%%%%%%%% Many physical systems contain conserved quantities: \begin{equation} I(X_t)=constant. \end{equation} During prediction, conservation error is measured as: \begin{equation} E_{cons} = \frac{1}{T} \sum_t | I(\hat{X}_t)-I(X_0) |. \end{equation} The evaluation includes: \begin{itemize} \item energy conservation; \item mass conservation; \item momentum preservation; \item system invariants. \end{itemize} %%%%%%%%%%%%%%%%%%%%%%%%%%%%%%%%%%%%%%%%%%%%%%%%%%%%%%%%%%%%%%%%%%%%%%%%%%%%%%%% \subsubsection{Lorenz-63 Physical Consistency} %%%%%%%%%%%%%%%%%%%%%%%%%%%%%%%%%%%%%%%%%%%%%%%%%%%%%%%%%%%%%%%%%%%%%%%%%%%%%%%% Although Lorenz-63 does not represent a conservation law system, it contains well-defined dynamical properties. The evaluation focuses on attractor preservation. The predicted trajectory: \begin{equation} \hat{X}(t) \end{equation} is compared with the reference attractor: \begin{equation} A_{true}. \end{equation} The attractor deviation is: \begin{equation} E_A = D( A_{true}, A_{pred} ). \end{equation} A physically consistent model should preserve: \begin{itemize} \item bounded chaotic evolution; \item attractor geometry; \item statistical trajectory properties. \end{itemize} %%%%%%%%%%%%%%%%%%%%%%%%%%%%%%%%%%%%%%%%%%%%%%%%%%%%%%%%%%%%%%%%%%%%%%%%%%%%%%%% \subsubsection{Reaction-Diffusion Physical Constraints} %%%%%%%%%%%%%%%%%%%%%%%%%%%%%%%%%%%%%%%%%%%%%%%%%%%%%%%%%%%%%%%%%%%%%%%%%%%%%%%% For reaction-diffusion systems, physical consistency is evaluated using the PDE residual: \begin{equation} R(u) = \frac{\partial u}{\partial t} - D\nabla^2u - R(u). \end{equation} The residual error is: \begin{equation} E_{PDE} = \frac{1}{N} \sum_i | R(u_i) |^2. \end{equation} PHYSGEN is expected to achieve lower residual values due to explicit integration of physical constraints. %%%%%%%%%%%%%%%%%%%%%%%%%%%%%%%%%%%%%%%%%%%%%%%%%%%%%%%%%%%%%%%%%%%%%%%%%%%%%%%% \subsubsection{Navier--Stokes Physical Constraints} %%%%%%%%%%%%%%%%%%%%%%%%%%%%%%%%%%%%%%%%%%%%%%%%%%%%%%%%%%%%%%%%%%%%%%%%%%%%%%%% For incompressible fluid dynamics, the main physical constraint is: \begin{equation} \nabla\cdot u=0. \end{equation} The divergence error is: \begin{equation} E_{div} = \frac{1}{N} \sum_i | \nabla\cdot\hat{u}_i |^2. \end{equation} Additionally, kinetic energy evolution is monitored: \begin{equation} E_k = \frac{1}{2} \int |\hat{u}|^2 dx. \end{equation} A physically consistent prediction should avoid artificial energy generation or dissipation. %%%%%%%%%%%%%%%%%%%%%%%%%%%%%%%%%%%%%%%%%%%%%%%%%%%%%%%%%%%%%%%%%%%%%%%%%%%%%%%% \subsubsection{Comparison with Baseline Models} %%%%%%%%%%%%%%%%%%%%%%%%%%%%%%%%%%%%%%%%%%%%%%%%%%%%%%%%%%%%%%%%%%%%%%%%%%%%%%%% The expected qualitative behavior is summarized in Table~\ref{tab:physics_consistency}. \begin{table}[h] \centering \caption{ Physical consistency comparison between PHYSGEN and baseline models. } \label{tab:physics_consistency} \begin{tabular}{lccc} \hline \textbf{Model} & \textbf{Prediction Error} & \textbf{Physics Violation} & \textbf{Long-Term Validity} \\ \hline MLP & Medium & High & Low \\ LSTM & Medium & High & Low \\ GNN & Low & Medium & Medium \\ GNS & Low & Medium & Medium \\ PINN & Medium & Low & Medium \\ FNO & Low & Medium & Medium \\ PHYSGEN & Low & Low & High \\ \hline \end{tabular} \end{table} %%%%%%%%%%%%%%%%%%%%%%%%%%%%%%%%%%%%%%%%%%%%%%%%%%%%%%%%%%%%%%%%%%%%%%%%%%%%%%%% \subsubsection{Physical Error Evolution} %%%%%%%%%%%%%%%%%%%%%%%%%%%%%%%%%%%%%%%%%%%%%%%%%%%%%%%%%%%%%%%%%%%%%%%%%%%%%%%% To analyze stability of physical constraints, the evolution of physical error is measured: \begin{equation} E_{phys}(k) = \| \mathcal{F} ( \hat{X}_{t+k} ) \|. \end{equation} The expected result is that PHYSGEN maintains bounded physical error during long rollout: \begin{equation} \lim_{k\rightarrow\infty} E_{phys}(k) < C. \end{equation} %%%%%%%%%%%%%%%%%%%%%%%%%%%%%%%%%%%%%%%%%%%%%%%%%%%%%%%%%%%%%%%%%%%%%%%%%%%%%%%% \subsubsection{Figure Representation} %%%%%%%%%%%%%%%%%%%%%%%%%%%%%%%%%%%%%%%%%%%%%%%%%%%%%%%%%%%%%%%%%%%%%%%%%%%%%%%% The corresponding visualization is: \begin{figure}[h] \centering \includegraphics[width=0.9\linewidth]{figures/physics_consistency.pdf} \caption{ Physical constraint violation during long-horizon prediction. Lower values indicate better preservation of physical laws. } \label{fig:physics_consistency} \end{figure} %%%%%%%%%%%%%%%%%%%%%%%%%%%%%%%%%%%%%%%%%%%%%%%%%%%%%%%%%%%%%%%%%%%%%%%%%%%%%%%% \subsubsection{Discussion} %%%%%%%%%%%%%%%%%%%%%%%%%%%%%%%%%%%%%%%%%%%%%%%%%%%%%%%%%%%%%%%%%%%%%%%%%%%%%%%% The physical consistency evaluation demonstrates the difference between accurate prediction and physically meaningful prediction. Purely data-driven approaches may approximate observed trajectories while violating underlying physical constraints. PHYSGEN addresses this limitation by integrating physical information directly into graph evolution, resulting in trajectories that remain closer to the physical state manifold. %%%%%%%%%%%%%%%%%%%%%%%%%%%%%%%%%%%%%%%%%%%%%%%%%%%%%%%%%%%%%%%%%%%%%%%%%%%%%%%% \subsubsection{Summary} %%%%%%%%%%%%%%%%%%%%%%%%%%%%%%%%%%%%%%%%%%%%%%%%%%%%%%%%%%%%%%%%%%%%%%%%%%%%%%%% The results of physical consistency evaluation demonstrate that PHYSGEN is designed not only as a forecasting model, but as a physically constrained dynamical system learner. By minimizing both prediction error and physical violations, PHYSGEN provides a more reliable framework for scientific simulation and digital twin applications. \] 
